\chapter{What was measured}
\label{ch:measured}

% Author's note: every number in this chapter is from the analysis notes
% of 2026-11-23 (campaign-2026-11d/cap-analysis.md, scripts listed in
% campaign-2026-11d/measurements/README.md), except the tiers of the
% literature search in sec:measured-tiers, which are from that search's
% report (campaign-2026-11d/lit-screening/work/report.md).  The scripts
% read private sessions and are not in this repository.

The defaults of \cref{ch:procedure} were set after four looks at
searches and screens of our own, all made on 2026-11-23.
Two are cases from one author's work: a review of one security question,
searched first with a limit and later without one
(\cref{sec:measured-limit}), and one book chapter's sixteen sessions,
screened in full (\cref{sec:measured-chapter}).
The third is a test of the classifier on fifty papers
(\cref{sec:measured-classifier}), the fourth the screening tiers of the
literature search of \cref{ch:lit-search} (\cref{sec:measured-tiers}).

What these measurements are matters for how far they carry.
They are a few cases, from one author's searches, with one language
model and one agent screener; they show what happened in these cases,
not what happens in general.
Where they say a work bears on a claim, that is the verdict of a
classifier or an agent, not of a person who checked it; where they say a
work was cited, that is what the authors chose from what they had found.
The scripts that computed them read the author's own sessions, which are
not published, so the numbers cannot be recomputed from this repository;
the literature search is the exception in part, since its search record
and the records it judged bearing are printed in the chapters that use
it.

\section{A limit on the search, and no limit}
\label{sec:measured-limit}

A review of one security question was searched with a limit of 40
results per query and database, written up, and later searched again
with the same queries and no limit.
We joined the two sets of exported records by paper.

\begin{description}
\item[Records] The search without a limit returned 8,267 records
  against 1,873 with it, 6,394 more.
\item[Bearing] Of the new records, 139 (2.2\,\%) were judged to bear on
  the question, against 7 to 12\,\% of the records the limited search
  had returned (the share differs by sub-question).
\item[By database] The new records came from OpenAlex (4,440, of which
  2.4\,\% bearing), IEEE Xplore (1,036; 0.4\,\%), Scopus (411; 5.4\,\%),
  arXiv (306; 0.3\,\%) and Web of Science (201; 3.5\,\%).
\item[Cited sources] The bibliography gained seven sources.
  Four were not among the limited search's records, two were (one of
  them discarded there), and one is in neither export.
  Of the cited sources found in the exports, 52 of 56 were already among
  the limited search's records.
\item[Conclusions] No result was overturned.
  One of the four sources the limit had hidden was the only review
  against one of the hypotheses.
  With the limit, the text had said that the search found no review
  against that hypothesis, which was true only of the limited search.
\end{description}

Two further changes to the review came from elsewhere: a study that no
query returned, found by comparing with another report, and four
sentences narrowed and two sources dropped after reading the cited works
in full.

\section{One book chapter's sixteen sessions}
\label{sec:measured-chapter}

A book chapter backs three claims with sixteen sessions, 36,808 records
in all, every one of them screened.
Before that screen the chapter had been written from a keyword pass,
without screening the records, followed by the check by meaning and the
repair queries it led to.

\paragraph{What the full screen added.}
The keyword pass gave 15 sources and the repair queries 4 counter
sources.
The finished chapter cites 16 works: 12 of the 15 and the 4 from the
repair queries.
The screen of all 36,808 records added no cited work.
The conclusion for each claim stayed as it was (one claim backed, one
narrower than first stated, one not backed); four works moved from
supporting to qualifying and one phrase was dropped, all after reading
the works in full.

\paragraph{Repeated screening.}
The 36,808 records were 20,678 distinct works.
Within one claim, 37\,\% of the screening was of a work already in
another session of the same claim (42\,\% and 34\,\% for the two claims
with several sessions), because each query had a session of its own.

\paragraph{A limit, simulated per list.}
From the cached answers of twelve queries, 68 lists in all, we took the
first $L$ of every list and counted what was left
(\cref{tab:measured-cap}).
Semantic Scholar's lists were ordered by citations, as its bulk search
orders a limited search.
The one cited work beyond place 500 was at place 512 of an OpenAlex
list.
Of the 16 cited works, 14 stood within the first 40 places of some
list, 13 of them in an OpenAlex or arXiv list and one in a Web of
Science list.

\begin{table}
  \begin{sidecaption}{%
      What the first $L$ places of each list would have held, of 19,607
      works in 68 lists of one book chapter, simulated from the cached
      answers.
      Works: the share of all works.
      Cited: of the chapter's 16 cited works.
      Bearing: the share of the works judged bearing.
    }[tab:measured-cap]
    \begin{tabular}{@{}rrrr@{}}
      \toprule
      $L$ & Works & Cited & Bearing \\
      \midrule
      40 & 1,551 (8\,\%) & 14 & 24\,\% \\
      100 & 16\,\% & 15 & 42\,\% \\
      200 & 27\,\% & 15 & 58\,\% \\
      500 & 48\,\% & 15 & 81\,\% \\
      1,000 & 66\,\% & 16 & 92\,\% \\
      \bottomrule
    \end{tabular}
  \end{sidecaption}
\end{table}

\paragraph{Per list or per question.}
For the two claims with several queries we took the lists of the
databases that order by relevance (OpenAlex, arXiv and IEEE Xplore;
14 lists from 5 queries and 19 from 7) and compared two ways to cut them.
The first $K$ places of every list, united, held 481 works and 8 of the
16 cited at $K=20$, 897 works and 14 cited at $K=40$, and 1,991 works
and 15 cited at $K=100$.
A fixed number $N$ of works per claim, taken in order of each work's best
place in any list, held 4 cited at $N=100$ (200 works for the two claims;
the 100th work had best place 10 and 7), 7 at 200, 10 at 300, 14 at
500 (1,000 works) and 15 at 1,000 (2,000 works).
So in this case a number per list reached as many cited works with
slightly fewer records than a number per claim (14 with 897 works against
1,000), and a small number per claim stopped near the top of every list:
100 per claim reached only the tenth place.

\paragraph{Where the bearing works sit.}
\Cref{tab:measured-density} gives the share of bearing works by place
in each database's lists.
In OpenAlex, arXiv and IEEE Xplore the share falls with the place; in
Scopus and Web of Science it does not, and those lists were ordered by
date: every neighbouring pair of Scopus records was in non-increasing
year order (mean year 2026.0 in the first quarter of a list, 2022.0 in
the last), and 89\,\% of Web of Science's pairs were, while OpenAlex,
arXiv and IEEE Xplore were not ordered by year.
Scholar sent no sort order to Scopus or Web of Science then; it now asks
both for relevance (\cref{sec:answer-orders}).

\begin{table}
  \begin{sidecaption}{%
      Share of works judged bearing, by place in each database's lists,
      in one book chapter's sessions.
      Semantic Scholar's lists were ordered by citations, Scopus's and
      Web of Science's by date.
    }[tab:measured-density]
    \begin{tabular}{@{}lrrrrr@{}}
      \toprule
      & 1--40 & 41--100 & 101--200 & 201--500 & 501-- \\
      \midrule
      OpenAlex & 42\,\% & 33\,\% & 29\,\% & 22\,\% & 11\,\% \\
      arXiv & 38\,\% & 24\,\% & 14\,\% & 10\,\% & 3\,\% \\
      IEEE Xplore & 23\,\% & 10\,\% & 8\,\% & 6\,\% & 2\,\% \\
      Semantic Scholar & 30\,\% & 26\,\% & 25\,\% & 23\,\% & 14\,\% \\
      Scopus & 26\,\% & 29\,\% & 29\,\% & 27\,\% & 18\,\% \\
      Web of Science & 24\,\% & 28\,\% & 28\,\% & 26\,\% & 34\,\% \\
      \bottomrule
    \end{tabular}
  \end{sidecaption}
\end{table}

\paragraph{What each database contributed.}
\Cref{tab:measured-providers} counts, per database, the works it
returned in the sixteen sessions, those no other database returned, the
bearing and the cited ones.
OpenAlex and arXiv together returned all 16 cited works.
The share of bearing works rose with the number of databases that
returned a work: 4\,\% for one, 8\,\% for two, 10\,\% for three and
19\,\% for four.

\begin{table}
  \begin{sidecaption}{%
      What each database contributed to one book chapter's sixteen
      sessions: works returned, works only it returned, works judged
      bearing, cited works, and cited works only it returned.
    }[tab:measured-providers]
    \begin{tabular}{@{}lrrrrr@{}}
      \toprule
      & Works & Only it & Bearing & Cited & Only it, cited \\
      \midrule
      OpenAlex & 13,450 & 6,345 & 1,146 & 14 & 4 \\
      Semantic Scholar & 10,191 & 2,858 & 850 & 10 & 0 \\
      arXiv & 5,019 & 1,757 & 259 & 8 & 0 \\
      Scopus & 3,537 & 496 & 540 & 5 & 0 \\
      Web of Science & 2,252 & 279 & 404 & 4 & 0 \\
      IEEE Xplore & 1,436 & 612 & 60 & 0 & 0 \\
      \bottomrule
    \end{tabular}
  \end{sidecaption}
\end{table}

\paragraph{Cost.}
A pilot classification in one of the sessions used about 158,000 tokens
of a Claude Sonnet model per 100 records.

\section{The classifier in batches}
\label{sec:measured-classifier}

We gave \texttt{scholar llm classify} the same 50 papers of one claim, all
with abstracts and no examples in the prompt, in three shapes: one batch
of 50, five batches of 10 (twice), and one paper per request.
The model was \texttt{anthropic/\allowbreak claude-sonnet-5-5}, and every request
was a real call to it (\cref{tab:measured-batches}).

\begin{table}
  \begin{sidecaption}{%
      One classifier, 50 papers, three batch sizes.
      Retried: requests whose answer had to be asked again.
      Answered: papers with a usable answer at the end.
      Seconds: the run's duration.
      Conf.: the mean confidence the model gave.
      Characters: of the prompts sent, retries included.
    }[tab:measured-batches]
    \footnotesize
    \begin{tabular}{@{}lrrrrrrr@{}}
      \toprule
      Batches & Requests & Retried & Seconds & Answered & Kept &
      Conf. & Characters \\
      \midrule
      1 of 50 & 1 & 0 & 57 & 50 & 26 & 0.74 & 104k \\
      5 of 10 & 5 & 1 & 91 & 50 & 27 & 0.78 & 131k \\
      5 of 10, again & 5 & 1 & 85 & 50 & 26 & 0.78 & 131k \\
      50 of 1 & 50 & 13 & 513 & 50 & 26 & 0.79 & 328k \\
      \bottomrule
    \end{tabular}
  \end{sidecaption}
\end{table}

No paper was skipped in any shape.
Every retry was for a first answer without a category tag, and after the
one retry every paper had a usable answer.
The two identical runs in batches of ten agreed on 49 of the 50
kept-or-discarded verdicts and on 46 of the categories; runs with
different batch sizes agreed on 45 to 48 verdicts and 38 to 43
categories.
Six papers changed status in at least one run, in no one direction: 26
or 27 papers were kept in every run.
Against the verdicts stored before (an earlier classifier run in batches
of 100 for 46 papers, an agent for 4; not a reference standard) the runs
agreed on 44 to 47 statuses and 38 to 43 categories.

\paragraph{A prompt with the papers last, withdrawn.}
A version of the prompt that put the papers after the instructions was
tried on the same papers.
In one batch of 50 the first answer had no category tag on any of the
50 papers, and each paper was asked again on its own: 138,101 tokens in
and 15,014 out, 153 seconds.
In five batches of 10, three needed no retry (7,200 to 9,500 tokens in)
and two lacked the category on 9 and on 10 of their papers (25,700 and
28,000 tokens in).
After the retries all 50 answers were usable, and they agreed with the
earlier runs on 43 to 46 statuses and 40 to 44 categories.
The prompt went back to its earlier order, and the retry became one
request holding only the papers whose answer was unusable
(\cref{sec:retry-one-paper}).

\paragraph{The revised code.}
With the earlier order and the new retry, one batch of 50 took one
request and no retry: 34,898 tokens in, 6,830 out, 51 seconds, all 50
answered, 28 kept.
Five batches of 10 took five requests and no retry: 39,410 tokens in and
5,845 out in all (7,200 to 9,400 in each), 12 to 14 seconds each, all 50
answered, 28 kept.
The revised runs agreed with the first test's batch of 50 on 48 of 50
statuses and 43 categories, with its batches of ten on 49 and 47, and
with each other on 48 and 43.
Counting all runs: with the earlier prompt order a retry was needed in
none of the batches of 50 and in 2 of 10 batches of 10 in the first test,
and in none of the 6 runs of the revised code; with the papers last, in
3 of 6 runs, each time for most or all papers of the batch.
A paper cost about 835 tokens, in and out together, in a batch of 50, and
about 905 in batches of 10.

\paragraph{Limits.}
Fifty papers, one claim, one model, and one repeat of one shape; the
first test did not record tokens or price.

\section{The tiers of the literature search}
\label{sec:measured-tiers}

The literature search of \cref{ch:lit-search} screened a first tier and
sampled the rest, so it measured the tier as it went.
Its four questions are those of \cref{ch:stopping,ch:abbreviated,%
ch:databases,ch:llm-screening}.
Bearing means judged to support or to qualify the question's claim, by
the one agent that screened the search.

\paragraph{The parts of the first tier.}
\Cref{tab:measured-tier-parts} gives each part's records and how many
bore on the question.
The check by meaning's works were the densest part, 48 to 73\,\%
bearing against 5 to 15\,\% for the first 100 places of the keyword
answers.
Records returned by several databases added little: beyond the first
100 places, 2 of the 49 records that three or more databases returned
for the first question bore on it, and a seeded sample of 60 records
returned by three or more databases, for each of the second and third
questions, held none (and one adjacent record each).
The tier was large because the answers were: 26 of the 35 OpenAlex
answers held more than 100 records.

\begin{table}
  \begin{sidecaption}{%
      The parts of the first tier in the literature search: records
      screened, and in parentheses those bearing on the question.
      First 100: places 1 to 100 of each query's OpenAlex and arXiv
      answers.
      Several databases: returned by at least three (first question) or
      five (the others) databases, beyond the first 100.
      Check: the check by meaning's works that were in the sessions.
      Known: named studies not in the other parts (--: none reported).
    }[tab:measured-tier-parts]
    \begin{tabular}{@{}lrrrr@{}}
      \toprule
      & First 100 & Several databases & Check & Known \\
      \midrule
      \Cref{ch:stopping} & 401 (49) & 49 (2) & 38 (22) & 1 \\
      \Cref{ch:abbreviated} & 675 (36) & 1 (0) & 17 (10) & 1 \\
      \Cref{ch:databases} & 996 (54) & 17 (0) & 21 (10) & -- \\
      \Cref{ch:llm-screening} & 645 (98) & 12 (2) & 45 (33) & -- \\
      \bottomrule
    \end{tabular}
  \end{sidecaption}
\end{table}

\paragraph{The samples of the rest.}
For each question, 60 records were drawn at random from the records left
unscreened and screened as the first tier was (\cref{sec:search-tiers}).
Three of the 240 bore on a question, all three on the first; in the first
tier 254 of 2,788 did (\cref{tab:measured-samples}).
Pooled, 2 of the 172 sampled records that one database returned were
bearing, 1 of 37 returned by two, none of 15 returned by three and none
of 16 returned by four or more.
None of the three would have entered the first tier by any of its
rules: none was in an OpenAlex or arXiv answer to the first question's
queries, none was returned by three or more databases, and none was
among the check by meaning's works.
Two were returned only by Web of Science or Scopus, whose answers were
then not ordered by relevance, and the third by both.

\begin{table}
  \begin{sidecaption}{%
      Records bearing on the question in the seeded sample of the
      unscreened records and in the first tier, with the 95\,\% modified
      Wilson interval of the sample's share
      (\cref{sec:llm-sample-interval}).
    }[tab:measured-samples]
    \begin{tabular}{@{}lllr@{}}
      \toprule
      & Sample & Interval & First tier \\
      \midrule
      \Cref{ch:stopping} & 3 of 60 (5\,\%) & 1.4 to 13.7\,\% & 53 of 454 \\
      \Cref{ch:abbreviated} & 0 of 60 & 0 to 6.0\,\% & 38 of 679 \\
      \Cref{ch:databases} & 0 of 60 & 0 to 6.0\,\% & 58 of 1,003 \\
      \Cref{ch:llm-screening} & 0 of 60 & 0 to 6.0\,\% & 105 of 652 \\
      All & 3 of 240 (1.25\,\%) & 0.3 to 3.6\,\% & 254 of 2,788 \\
      \bottomrule
    \end{tabular}
  \end{sidecaption}
\end{table}
